
\section{Pair production in particle collisions}

\SourcePage{491}{496}

The production of an electron pair in the collision of two charged particles is described by diagrams of two types:

% [Feynman diagrams (100.1): a) two-photon diagram with upper lines for colliding particles and lower line for produced pair e⁻e⁺; b) one-photon exchange diagram with pair produced from field of second particle]

The two upper solid lines correspond to the colliding particles, the lower one to the pair being produced.

Let us consider in the ultrarelativistic case the collision of two heavy particles (nuclei). The change in the state of motion of the particles themselves in such a collision can be neglected, i.e.\ they can be considered as sources of an external field\footnote{The case of the collision of two light particles (electrons), the change in the motion of which cannot be neglected, is considerably more complicated. See about this the book by V.~Ya.~Baier, V.~M.~Katkov and V.~S.~Fadin indicated on p.~454.}. To this there correspond two diagrams of the first type:

% [Feynman diagrams (100.2): two diagrams showing pair production from two-photon exchange, with momentum labels q⁽¹⁾, q⁽²⁾ for Fourier components of fields of two particles]

where $q^{(1)}$, $q^{(2)}$ are the ``momenta'' of the Fourier components of the fields of the two particles.

The potential $A^\mu = (A_0,\,\mathbf{A})$, created by a uniformly moving with velocity $\mathbf{v}$ classical particle, satisfies the equations
\[
\Box A_0 = -4\pi Ze\,\delta(\mathbf{r} - \mathbf{v}t - \mathbf{r}_0), \quad \Box\mathbf{A} = -4\pi Ze\mathbf{v}\,\delta(\mathbf{r} - \mathbf{v}t - \mathbf{r}_0).
\]

Its Fourier components:
\[
A_0(\omega,\mathbf{k}) = -\frac{8\pi^2 Ze}{\omega^2 - \mathbf{k}^2}\,e^{-i\mathbf{k}\mathbf{r}_0}\,\delta(\omega - \mathbf{k}\mathbf{v})
\]
and analogously for $\mathbf{A}(\omega,\mathbf{k})$. In four-dimensional form
\[
A^\mu(q) = -\frac{8\pi^2 Ze}{q^2}\,e^{iqx_0}\,U^\mu\,\delta(Uq),
\]
where $U$ is the 4-velocity of the particle, and the 4-vector $x_0 = (0,\,\mathbf{r}_0)$. If nucleus 1 is at rest at the origin ($\mathbf{r}_0^{(1)} = 0$), then $\boldsymbol{\rho} \equiv \mathbf{r}_0^{(2)}$ is the impact parameter vector (in the plane perpendicular to
